\subsection{固定时域的闭环对照与边界状态}
下表每项均来自本地完整334日轨迹，统一真实初末6000 kWh，保持预报权重1、28日窗口、7条等权场景、321库存网格与三箱Markov反馈。第二个午夜视域随重规划时刻取48、42、36、30小时，固定视域则分别取48或72小时，评价终点统一截断。
\begin{table}[H]\centering\small
\caption{同配置下的时域诊断（万元）}
\begin{tabular}{lrrrr}\toprule
机制 & 第二个午夜 & 固定48小时 & 固定72小时 & 48至72小时节省\\\midrule
问题二 & 1384.42 & 1384.36 & 1381.79 & +2.57 \\
问题三 & 1346.18 & 1346.13 & 1347.26 & -1.13 \\
问题四日前 & 1465.46 & 1465.43 & 1459.31 & +6.12 \\
问题四日内 & 1425.21 & 1425.21 & 1424.40 & +0.81 \\
\bottomrule\end{tabular}\end{table}
改变时域还会改变完整成熟残差的可用集合和末值定价时段，因此该表识别配置变化的闭环总效应。固定72小时在6、12、18时起算时，规划末点落在日内；完整继续状态应包含未交付合同，而当前实现用库存线性末值近似。实际年度账单覆盖全部交付，这一近似影响规划继续价值。上述时域对照同时改变历史池和末值时段，不能独立量化边界近似误差；其识别需要在相同72小时条件下增加完整合同继续状态的匹配回放。主策略的滚动末点位于午夜。

问题二固定72小时相对两午夜视域，334日节省26,225.00元，日均78.52元。3、7、14日块95\%区间分别为$[-122.81,295.64]$、$[-101.14,288.44]$、$[-74.14,285.42]$元/日；节省日156、增支日177、近零日1，正累计月份5个。区间跨零表示该年度增量证据尚不足以确认稳定正收益，不表示真实效果已被证明为零。该诊断对应固定72小时候选，不替代正文的强跨日贪心对照。

\subsection{预报融合的二阶矩推导与开发诊断}
记官方和历史预测误差为$e_O,e_H$，融合误差为$e_\beta=\beta e_O+(1-\beta)e_H$。假设二阶矩有限，令
\begin{equation}
 A=\E[e_O^2],\quad B=\E[e_H^2],\quad C=\E[e_Oe_H],\quad D=A+B-2C.
\end{equation}
则
\begin{equation}
 \E[e_\beta^2]=B+2(C-B)\beta+D\beta^2,
 \qquad D=\E[(e_O-e_H)^2]\ge0.
\end{equation}
当$D>0$时，凸二次式在$[0,1]$的最小点为
\begin{equation}
 \beta^*=\operatorname{proj}_{[0,1]}\frac{B-C}{D}.
\end{equation}
证明由求导及闭区间投影直接得到。当$D=0$时，两误差几乎处处相等，所有权重具有相同MSE。二阶矩包含偏差平方，不能在有偏预测下直接以方差替代。将期望替换为同一组配对段平均，可得到对应样本MSE的最优权重。

\begin{table}[H]\centering\small\setlength{\tabcolsep}{3pt}
\caption{配对光伏误差矩与样本最优权重（二阶矩单位kW$^2$）}
\begin{tabular}{lrrrrrr}\toprule
一月窗口 & 段数 & $A$ & $B$ & $C$ & $D$ & $\widehat\beta^*$\\\midrule
15至24日 & 1440 & 10849.55 & 33705.63 & 5172.59 & 34210.01 & 0.8341 \\
25至31日 & 1008 & 23689.77 & 47340.10 & 12923.33 & 45183.21 & 0.7617 \\
\bottomrule\end{tabular}\end{table}
每个0、6、12、18时原点预测后续36段，每个目标在所属窗口计入一次。15至24日的离散候选权重为0、0.25、0.5、0.75、1，MAE分别为96.68、75.89、59.42、50.50、52.04 kW；因此0.75的选取具有该候选集的开发期MAE依据。25至31日中0.75的MSE为21,130.50 kW$^2$，样本最优权重0.7617对应21,124.30 kW$^2$，两者接近。

历史预测族已用1月15至31日选型，因此两个窗口均属于内部开发诊断；按原点截断未来输入不能消除预测族选型的日期重叠。这里不使用2至12月样本估计融合矩，也不将样本MSE最优权重解释为无偏留出估计或账单最优权重。全年账单受净需求尾部、合同调整费与库存约束共同影响，预测误差降低的经济作用需由固定时域下的完整配对执行验证。
